\section{Related Work}

\subsection{CTC Decoding and N-best Generation}

CTC \cite{graves2006ctc} trains an acoustic model to assign probability to token sequences without frame-level alignment supervision. The model outputs a distribution $p_t(k | x)$ over the vocabulary at each encoder frame $t$; the probability of a token sequence $y$ is obtained by marginalizing over all valid frame-to-token alignments via the forward-backward algorithm, which runs in $O(T \cdot |y|)$ time. Each frame-level prediction depends on the full acoustic context through the encoder but not on adjacent output decisions. This conditional independence at the output projection enables parallel decoding and makes CTC well-suited to latency-sensitive settings.

Greedy decoding selects the argmax label at each frame and collapses consecutive duplicates and blanks. It is fast and exact in the MAP sense over frame-level decisions, but it does not maximize the true CTC sequence probability $P_\text{CTC}(y|x)$: the per-frame argmax labels do not necessarily compose into the highest-probability token sequence after alignment marginalization. Prefix beam search \cite{graves2014towards} corrects this by maintaining a set of partial hypotheses and accumulating CTC probabilities over all valid alignments for each prefix. Prefix beam search produces a scored N-best list ordered by accumulated CTC probability and remains the standard CTC decoding algorithm.

A single beam search does not guarantee diverse candidates. When the model is confident, the top-$G$ hypotheses from a fixed beam may differ only in a few tokens, and the oracle WER of the N-best list barely undercuts the greedy WER. The k2 framework \cite{li2022k2} addresses this through lattice-based sampling: a first-pass lattice is built over a wide beam, and $G$ hypotheses are drawn from it using a temperature parameter (nbest\_scale) that controls arc-weight sharpness. At nbest\_scale=1.0, the oracle WER at $G=128$ on LibriSpeech dev-other reaches 3.53\%, against a greedy baseline of 6.02\% — a 2.49~pp oracle gap. Reducing nbest\_scale to 0.5 destroys approximately 90\% of this gap by collapsing the lattice distribution prematurely; all experiments in this work use nbest\_scale=1.0.

The primary acoustic model throughout this paper is a Zipformer-S \cite{yao2024zipformer} trained with the CR-CTC objective \cite{huang2024crctc}. CR-CTC augments the CTC loss with a consistency regularization term: the KL divergence between the output posteriors of two stochastically augmented views of the same utterance is added as a penalty. The effect on the alignment posterior $\gamma_t$ is measurable. On 100 dev-other utterances, the CR-CTC model produces a near-equal three-way split — dead frames (blank-dominated, $\gamma_t(\text{blank}) > 0.99$) at 34.2\%, active frames at 34.7\%, and ambiguous frames at 31.1\%. Unregularized CTC models concentrate 70–90\% of frames in the dead category \cite{graves2006ctc}; CR-CTC roughly halves that fraction by spreading alignment mass more evenly across active frames.

CTC posterior scores alone cannot close this gap: CTC scoring tracks acoustic confidence rather than linguistic plausibility, and the ranking signal degrades as the beam grows. Section~\ref{sec:exhaustion} verifies this across eleven CTC-internal scoring strategies, all of which fail to produce statistically significant WER improvements. An external linguistic signal is required.



\subsection{Language Model Rescoring for ASR}

Two strategies dominate LM integration in end-to-end ASR: shallow fusion and N-best rescoring. Shallow fusion \cite{gulcehre2015fusion} combines the acoustic model log-probability with an autoregressive LM log-probability at each decoding step:
\begin{equation*}
\log \tilde{P}(y_t | y_{<t}, x) \;=\; \log P_\text{AM}(y_t | x) \;+\; \lambda\, \log P_\text{LM}(y_t | y_{<t}).
\end{equation*}
The interpolation weight $\lambda$ is tuned on a held-out set. \citet{hori2017advances} apply this approach to attention-based encoder-decoder ASR and demonstrate consistent WER reductions across languages and model sizes. The strategy is straightforward to implement, but its limitation is that the LM must be queried at every beam expansion step, coupling inference latency to LM size.

N-best rescoring decouples acoustic decoding from LM scoring. The acoustic decoder runs first and produces a fixed candidate list; the LM scores each candidate offline, and the combined score determines the final selection. \citet{mikolov2010recurrent} show that recurrent neural network LMs substantially outperform n-gram models in this setting. Their work establishes the second-pass neural rescoring paradigm that modern systems still follow. \citet{shin2019effective} bring this to end-to-end ASR: a pre-trained transformer LM scoring the CTC N-best list achieves WER reductions comparable to shallow fusion without modifying the acoustic decoder architecture.

For CTC systems, WFST-based decoding \cite{miao2015eesen} offers a third option. The CTC decoding graph is composed with a compiled n-gram language model in a weighted finite-state transducer; the result enables joint search over acoustic and language model scores. Within the n-gram model class, this approach is exact. Its structural limitation is that the grammar must be compiled at a fixed vocabulary and n-gram order, so on-the-fly updates and neural LM integration are impractical.

All three approaches share a common form for the LM score: an autoregressive or n-gram conditional probability $P(y_t | y_{<t})$. Bidirectional masked LMs such as RoBERTa \cite{liu2019roberta} produce non-autoregressive posteriors that do not fit this interface and cannot be used in shallow fusion or WFST integration. \citet{salazar2020pll} introduce pseudo-log-likelihood (PLL) as a sentence-level score for masked LMs: each token position is individually masked, and the sum of per-position log-probabilities under the masked contexts defines the PLL for the complete sentence. PLL correlates with human grammaticality judgments and outperforms autoregressive LM scores on acceptability benchmarks \cite{salazar2020pll}. In prior ASR work, however, PLL serves only as a drop-in replacement for conventional LM scores in a rescoring-and-argmax pipeline.

This work changes the decoding objective rather than the scoring function. The PLL score defines a probability distribution over the N-best candidates, and MBR decoding then selects the hypothesis minimizing expected CER under that distribution — not the hypothesis with the highest score. When the CTC posterior concentrates mass around near-duplicate candidates, the argmax and MBR solutions diverge; the MBR solution tends to be better in these cases, as the Spearman-$\rho$ divergence analysis in §5.3 quantifies. To our knowledge, the combination of masked-LM posteriors with an MBR objective applied to CTC N-best lists has not been explored before this work.



\subsection{Minimum Bayes Risk Decoding}

MBR decoding selects the output that minimizes expected loss under a posterior over the hypothesis space:
\begin{equation*}
\hat{y}_\text{MBR} \;=\; \operatorname{argmin}_{y \in \mathcal{H}} \sum_{y' \in \mathcal{H}} Q(y') \, L(y, y'),
\end{equation*}
where $\mathcal{H}$ is the candidate set, $Q(y')$ is a probability distribution over hypotheses, and $L$ is a task-specific loss function. MAP decoding selects the hypothesis assigned highest model probability, without accounting for the cost structure $L$ — that is, without considering how wrong each candidate would be if each alternative turned out to be correct. \citet{kumar2004mbr} introduce the MBR formulation for machine translation and show that sentence-level BLEU MBR outperforms MAP decoding on multiple benchmarks; the gain is largest on outputs where several near-equivalent hypotheses share high model probability.

In machine translation, MBR has become a standard technique. \citet{eikema2020map} argue theoretically that MAP decoding is ill-suited to neural sequence generation: the probability mass of a well-trained model spreads across many paraphrase variants of the same meaning, so any single MAP hypothesis is unrepresentative of the distribution. Sampling-based MBR, which draws a large hypothesis set from the model and applies the MBR objective over the sample, produces more fluent and accurate translations. \citet{freitag2022high} scale this approach using learned neural metrics (BLEURT, COMET) as the utility function $L$ and report substantial improvements over beam search on WMT benchmarks. \citet{finkelstein2024mbr} extend the framework to training: NMT models fine-tuned to minimize expected loss produce better translations than cross-entropy-trained models; the result shows that MBR and the training objective can be unified.

MBR for ASR has precedent, though it predates the neural era. \citet{goel2000mbr} apply a minimum expected WER criterion to acoustic lattice rescoring in hybrid HMM-based systems: the posterior $Q(y')$ is estimated from the acoustic lattice, and the hypothesis minimizing expected WER under this posterior is selected in place of the MAP hypothesis. These results show consistent improvements over MAP decoding on the NIST Switchboard benchmark. In that work, however, $Q(y')$ is derived entirely from the acoustic model. No external linguistic signal enters the posterior; the lattice already encodes all available information.

When $Q(y')$ is the CTC acoustic distribution, MBR over the N-best list can only reweight candidates within what the acoustic model already knows. The Spearman rank correlation between CTC hypothesis scores and per-utterance WER degrades from $-0.347$ at $G=16$ to $-0.270$ at $G=128$ as near-duplicate candidates dilute the acoustic ranking signal — a 22\% drop. PLL degrades much more gracefully over the same range ($-0.484$ to $-0.461$, a 5\% drop), and the two signals diverge most sharply on utterances where the correct transcription differs substantially from the CTC argmax. Augmenting the posterior with PLL changes which candidates MBR selects in precisely these recoverable cases. The closest precedent is \citet{finkelstein2024mbr} in NMT, which uses neural metrics as the utility function $L$ while retaining an autoregressive model posterior; this paper changes the posterior $Q$ itself and operates in the ASR domain.



\subsection{RL and Sequence-Level Training for ASR}

The connection between policy gradient RL and sequence-level ASR training is well established. \citet{williams1992} introduces REINFORCE: the gradient of expected reward with respect to policy parameters equals the expected product of the reward and the gradient of the log-policy probability. Applied to sequence generation, the action is the output token sequence, the reward is a sequence-level metric such as WER, and the policy is the acoustic model distribution. \citet{prabhavalkar2018mwer} derive the MWER training loss directly from this framing:
\begin{equation*}
\mathcal{L}_\text{MWER} \;=\; \sum_{i=1}^{G} \hat{A}_i \;\log P_\text{CTC}(y_i|x),
\end{equation*}
where $\hat{A}_i = \text{WER}(y_i, y^*) - \frac{1}{G}\sum_{j} \text{WER}(y_j, y^*)$ is the mean-centered WER advantage over the N-best list. \citet{shannon2017} develop the same functional form and show that sequence-level fine-tuning after cross-entropy pretraining consistently reduces WER. Both works treat the REINFORCE identification as motivation rather than as an object of analysis.

REINFORCE gradient estimates have high variance, which motivates several reduction strategies. \citet{rennie2017self} propose self-critical sequence training (SCST): the reward baseline is the performance of the model's own greedy output at test time. This adaptive baseline generalizes the mean-reward baseline of \citet{williams1992} by conditioning on the actual test-time decoding decision. The Rao-Blackwell theorem \cite{casella1996rao} provides a more structural guarantee: replacing a random variable with its conditional expectation given a sufficient statistic produces an estimator with weakly lower variance. \citet{ranganath2014bbvi} apply Rao-Blackwellization to black-box variational inference by conditioning score-function gradient estimators on subsets of the latent variables. This reduces per-sample variance without changing the expected gradient.

Recent work on large language models revisits these ideas with importance-ratio clipping. GRPO \cite{shao2024grpo} extends REINFORCE with a group-relative advantage baseline and a PPO-style clipping term, improving training stability for LLM fine-tuning on reasoning benchmarks. The structure is the same as MWER: a finite action space, a non-differentiable reward, and a policy that is already near-converged before RL training begins. Section~\ref{sec:training} applies clipping of this type to the CR-CTC setting; it reduces the magnitude of degradation but does not reverse its direction.

The theoretical contribution of this work connects these ideas through the CTC backward pass. Standard MWER training computes the gradient of $\log P_\text{CTC}(y_i|x)$ via the CTC forward-backward algorithm. This returns the alignment posterior $\gamma_t(k|y_i)$ as the per-frame gradient weight. We show in Section~\ref{sec:theory} that this marginalization over alignments is exactly a Rao-Blackwellization of a single-alignment REINFORCE estimator: $\gamma_t(k|y_i)$ is the conditional expectation of the Viterbi or sampled per-frame indicator given the token sequence $y_i$, and the Rao-Blackwell theorem guarantees its variance is weakly lower. Empirically, the CTC-marginal gradient has 2.96$\times$ lower variance than the Viterbi estimator and 3.87$\times$ lower variance than a sampled single-path estimator. Any practitioner running standard MWER training via the CTC backward pass already obtains this variance reduction; the contribution of this paper is identifying and verifying the mechanism. The RL framing is limited to training-time optimization: at decode time, the model parameters are fixed, and the problem becomes one of action selection under uncertainty — Bayesian decision theory, developed in Section~\ref{sec:theory}.



